\section{Syntomic morphisms}
\label{section-syntomic}

\noindent
Syntomic ring maps are flat finitely presented ring maps all of whose fibers
are local complete intersections. We discuss general local complete
intersection ring maps in
More on Algebra, Section \ref{more-algebra-section-lci}.

\begin{definition}
\label{definition-lci}
A ring map $R \to S$ is called {\it syntomic}, or we say $S$ is a
{\it flat local complete intersection over $R$}
if it is flat, of finite presentation, and if all of its fibre rings
$S \otimes_R \kappa(\mathfrak p)$ are local complete intersections,
see Definition \ref{definition-lci-field}.
\end{definition}

\noindent
Clearly, an algebra over a field is syntomic over the field
if and only if it is a local complete intersection. Here is
a pleasing feature of this definition.

\begin{lemma}
\label{lemma-syntomic-descends}
\begin{slogan}
Being syntomic is fpqc local on the base.
\end{slogan}
Let $R \to S$ be a ring map.
Let $R \to R'$ be a faithfully flat ring map.
Set $S' = R'\otimes_R S$.
Then $R \to S$ is syntomic if and only if $R' \to S'$ is syntomic.
\end{lemma}

\begin{proof}
By Lemma \ref{lemma-finite-presentation-descends} and
Lemma \ref{lemma-flatness-descends} this holds for the property
of being flat and for the property of being of finite presentation.
The map $\Spec(R') \to \Spec(R)$ is surjective,
see Lemma \ref{lemma-ff-rings}. Thus it suffices to show
given primes $\mathfrak p' \subset R'$ lying over $\mathfrak p \subset R$
that $S \otimes_R \kappa(\mathfrak p)$ is a local complete
intersection if and only if $S' \otimes_{R'} \kappa(\mathfrak p')$
is a local complete intersection. Note that
$S' \otimes_{R'} \kappa(\mathfrak p') =
S \otimes_R \kappa(\mathfrak p)
\otimes_{\kappa(\mathfrak p)} \kappa(\mathfrak p')$.
Thus Lemma \ref{lemma-lci-field-change} applies.
\end{proof}

\begin{lemma}
\label{lemma-base-change-syntomic}
Any base change of a syntomic map is syntomic.
\end{lemma}

\begin{proof}
This is true for being flat, for being of finite presentation,
and for having local complete intersections as fibres by
Lemmas \ref{lemma-flat-base-change}, \ref{lemma-compose-finite-type} and
\ref{lemma-lci-field-change}.
\end{proof}

\begin{lemma}
\label{lemma-local-syntomic}
Let $R \to S$ be a ring map.
Suppose we have $g_1, \ldots g_m \in S$ which generate the
unit ideal such that each $R \to S_{g_i}$ is syntomic.
Then $R \to S$ is syntomic.
\end{lemma}

\begin{proof}
This is true for being flat and for being of finite presentation by
Lemmas \ref{lemma-flat-localization} and \ref{lemma-cover-upstairs}.
The property of having fibre rings which are local complete intersections
is local on $S$ by its very definition, see
Definition \ref{definition-lci-field}.
\end{proof}

\begin{definition}
\label{definition-relative-global-complete-intersection}
Let $R \to S$ be a ring map. We say that $R \to S$ is
a {\it relative global complete intersection} if there exists
a presentation $S = R[x_1, \ldots, x_n]/(f_1, \ldots, f_c)$ and
every nonempty fibre of $\Spec(S) \to \Spec(R)$ has dimension $n - c$.
We will say ``let $S = R[x_1, \ldots, x_n]/(f_1, \ldots, f_c)$ be a relative
global complete intersection'' to indicate this situation.
\end{definition}

\noindent
The following lemma is occasionally useful to find
global presentations.

\begin{lemma}
\label{lemma-huber}
Let $S$ be a finitely presented $R$-algebra which has a presentation
$S = R[x_1, \ldots, x_n]/I$ such that $I/I^2$ is free over $S$. Then
$S$ has a presentation $S = R[y_1, \ldots, y_m]/(f_1, \ldots, f_c)$
such that $(f_1, \ldots, f_c)/(f_1, \ldots, f_c)^2$ is free with
basis given by the classes of $f_1, \ldots, f_c$.
\end{lemma}

\begin{proof}
Note that $I$ is a finitely generated ideal by
Lemma \ref{lemma-finite-presentation-independent}.
Let $f_1, \ldots, f_c \in I$ be elements which map to a basis of $I/I^2$.
By Nakayama's lemma (Lemma \ref{lemma-NAK})
there exists a $g \in 1 + I$ such that
$$
g \cdot I \subset (f_1, \ldots, f_c)
$$
and $I_g \cong (f_1, \ldots, f_c)_g$. Hence we see that
$$
S \cong R[x_1, \ldots, x_n]/(f_1, \ldots, f_c)[1/g]
\cong R[x_1, \ldots, x_n, x_{n + 1}]/(f_1, \ldots, f_c, gx_{n + 1} - 1)
$$
as desired. It follows that $f_1, \ldots, f_c,gx_{n + 1} - 1$
form a basis for
$(f_1, \ldots, f_c, gx_{n + 1} - 1)/(f_1, \ldots, f_c, gx_{n + 1} - 1)^2$
for example by applying Lemma \ref{lemma-principal-localization-NL}.
\end{proof}



\begin{example}
\label{example-factor-polynomials}
Let $n , m \geq 1$ be integers. Consider the ring map
\begin{eqnarray*}
R = \mathbf{Z}[a_1, \ldots, a_{n + m}]
& \longrightarrow &
S = \mathbf{Z}[b_1, \ldots, b_n, c_1, \ldots, c_m] \\
a_1 & \longmapsto & b_1 + c_1 \\
a_2 & \longmapsto & b_2 + b_1 c_1 + c_2 \\
\ldots & \ldots & \ldots \\
a_{n + m} & \longmapsto & b_n c_m
\end{eqnarray*}
In other words, this is the unique ring map of polynomial rings
as indicated such that the polynomial factorization
$$
x^{n + m} + a_1 x^{n + m - 1} + \ldots + a_{n + m}
=
(x^n + b_1 x^{n - 1} + \ldots + b_n)
(x^m + c_1 x^{m - 1} + \ldots + c_m)
$$
holds. Note that $S$ is generated by $n + m$ elements over $R$
(namely, $b_i, c_j$) and that there are $n + m$ equations
(namely $a_k = a_k(b_i, c_j)$). In order to show that
$S$ is a relative global complete intersection over $R$ it suffices
to prove that all fibres have dimension $0$.

\medskip\noindent
To prove this, let $R \to k$ be a
ring map into a field $k$. Say $a_i$ maps to $\alpha_i \in k$.
Consider the fibre ring $S_k = k \otimes_R S$. Let $k \to K$ be
a field extension. A $k$-algebra map of $S_k \to K$ is the same thing as
finding $\beta_1, \ldots, \beta_n, \gamma_1, \ldots, \gamma_m \in K$
such that
$$
x^{n + m} + \alpha_1 x^{n + m - 1} + \ldots + \alpha_{n + m}
=
(x^n + \beta_1 x^{n - 1} + \ldots + \beta_n)
(x^m + \gamma_1 x^{m - 1} + \ldots + \gamma_m).
$$
Hence we see there are at most finitely many choices of
such $n + m$-tuples in $K$. This proves that all fibres
have finitely many closed points (use Hilbert's Nullstellensatz
to see they all correspond to solutions in $\overline{k}$ for example)
and hence that $R \to S$ is a relative global complete intersection.

\medskip\noindent
Another way to argue this is to show
$\mathbf{Z}[a_1, \ldots, a_{n + m}] \to
\mathbf{Z}[b_1, \ldots, b_n, c_1, \ldots, c_m]$ is actually
also a {\it finite} ring map. Namely, by Lemma \ref{lemma-polynomials-divide}
each of $b_i, c_j$ is integral over $R$, and hence $R \to S$ is
finite by Lemma \ref{lemma-characterize-integral}.
\end{example}

\begin{example}
\label{example-roots-universal-polynomial}
Consider the ring map
\begin{eqnarray*}
R = \mathbf{Z}[a_1, \ldots, a_n]
& \longrightarrow &
S = \mathbf{Z}[\alpha_1, \ldots, \alpha_n] \\
a_1 & \longmapsto &
\alpha_1 + \ldots + \alpha_n \\
\ldots & \ldots & \ldots \\
a_n & \longmapsto & \alpha_1 \ldots \alpha_n
\end{eqnarray*}
In other words this is the unique ring map of polynomial
rings as indicated such that
$$
x^n + a_1 x^{n - 1} + \ldots + a_n
=
\prod\nolimits_{i = 1}^n (x + \alpha_i)
$$
holds in $\mathbf{Z}[\alpha_i, x]$. Another way to say this
is that $a_i$ maps to the $i$th elementary symmetric function
in $\alpha_1, \ldots, \alpha_n$. By the usual theory of elementary
symmetric polynomials (details omitted) the ring $S$ is finite
free over $R$ with basis the elements
$\alpha_1^{e_1} \alpha_2^{e_2} \ldots \alpha_n^{e_n}$ with
$0 \leq e_i \leq n - i$. A fortiori, the fibre rings of
$R \to S$ are finite and hence have dimension $0$.
On the other hand, $S$ is generated by
$n$ elements over $R$ subject to $n$ equations. Hence $S$
is a relative global complete intersection over $R$.
Since the rank of $S$ over $R$ is positive, we also see
that $S$ is faithfully flat over $R$.
\end{example}

\begin{lemma}
\label{lemma-base-change-relative-global-complete-intersection}
Let $S = R[x_1, \ldots, x_n]/(f_1, \ldots, f_c)$ be a
relative global complete intersection
(Definition \ref{definition-relative-global-complete-intersection})
\begin{enumerate}
\item For any $R \to R'$ the base change
$R' \otimes_R S = R'[x_1, \ldots, x_n]/(f_1, \ldots, f_c)$ is a relative
global complete intersection.
\item For any $g \in S$ which is the image of $h \in R[x_1, \ldots, x_n]$
the ring
$S_g = R[x_1, \ldots, x_n, x_{n + 1}]/(f_1, \ldots, f_c, hx_{n + 1} - 1)$
is a relative global complete intersection.
\item If $R \to S$ factors as $R \to R_f \to S$ for some $f \in R$.
Then the ring $S = R_f[x_1, \ldots, x_n]/(f_1, \ldots, f_c)$
is a relative global complete intersection over $R_f$.
\end{enumerate}
\end{lemma}

\begin{proof}
By Lemma \ref{lemma-dimension-preserved-field-extension}
the fibres of a base change have the same dimension as the
fibres of the original map. Moreover
$R' \otimes_R R[x_1, \ldots, x_n]/(f_1, \ldots, f_c)
= R'[x_1, \ldots, x_n]/(f_1, \ldots, f_c)$. Thus (1) follows.
The proof of (2) is that
the localization at one element can be described as
$S_g \cong S[x_{n + 1}]/(gx_{n + 1} - 1)$.
Assertion (3) follows from (1) since under the assumptions of (3) we have
$R_f \otimes_R S \cong S$.
\end{proof}

\begin{lemma}
\label{lemma-localize-relative-complete-intersection}
Let $R$ be a ring. Let $S = R[x_1, \ldots, x_n]/(f_1, \ldots, f_c)$.
We will find $h \in R[x_1, \ldots, x_n]$ which maps to $g \in S$
such that
$$
S_g = R[x_1, \ldots, x_n, x_{n + 1}]/(f_1, \ldots, f_c, hx_{n + 1} - 1)
$$
is a relative global complete intersection with a presentation as in
Definition \ref{definition-relative-global-complete-intersection}
in each of the following cases:
\begin{enumerate}
\item Let $I \subset R$ be an ideal. If the fibres of
$\Spec(S/IS) \to \Spec(R/I)$ have dimension $n - c$, then we can
find $(h, g)$ as above such that $g$ maps to $1 \in S/IS$.
\item Let $\mathfrak p \subset R$ be a prime. If
$\dim(S \otimes_R \kappa(\mathfrak p)) = n - c$, then we can
find $(h, g)$ as above such that $g$ maps to a unit of
$S \otimes_R \kappa(\mathfrak p)$.
\item Let $\mathfrak q \subset S$ be a prime lying over
$\mathfrak p \subset R$. If $\dim_{\mathfrak q}(S/R) = n - c$, then we can
find $(h, g)$ as above such that $g \not \in \mathfrak q$.
\end{enumerate}
\end{lemma}

\begin{proof}
Ad (1). By Lemma \ref{lemma-dimension-fibres-bounded-open-upstairs}
there exists an open subset $W \subset \Spec(S)$ containing $V(IS)$
such that all fibres of $W \to \Spec(R)$ have dimension $\leq n - c$.
Say $W = \Spec(S) \setminus V(J)$. Then $V(J) \cap V(IS) = \emptyset$
hence we can find a $g \in J$ which maps to $1 \in S/IS$.
Let $h \in R[x_1, \ldots, x_n]$ be any preimage of $g$.

\medskip\noindent
Ad (2). By Lemma \ref{lemma-dimension-fibres-bounded-open-upstairs}
there exists an open subset $W \subset \Spec(S)$ containing
$\Spec(S \otimes_R \kappa(\mathfrak p))$
such that all fibres of $W \to \Spec(R)$ have dimension $\leq n - c$.
Say $W = \Spec(S) \setminus V(J)$. Then
$V(J \cdot S \otimes_R \kappa(\mathfrak p)) = \emptyset$.
Hence we can find a $g \in J$ which maps to a unit in
$S \otimes_R \kappa(\mathfrak p)$ (details omitted).
Let $h \in R[x_1, \ldots, x_n]$ be any preimage of $g$.

\medskip\noindent
Ad (3). By Lemma \ref{lemma-dimension-fibres-bounded-open-upstairs}
there exists a $g \in S$, $g \not \in \mathfrak q$
such that all nonempty fibres of $R \to S_g$
have dimension $\leq n - c$. Let $h \in R[x_1, \ldots, x_n]$
be any element that maps to $g$.
\end{proof}

\noindent
The following lemma says we can do absolute Noetherian
approximation for relative global complete intersections.

\begin{lemma}
\label{lemma-relative-global-complete-intersection-Noetherian}
Let $R$ be a ring. Let $S = R[x_1, \ldots, x_n]/(f_1, \ldots, f_c)$
be a relative global complete intersection
(Definition \ref{definition-relative-global-complete-intersection}).
There exist a finite type $\mathbf{Z}$-subalgebra $R_0 \subset R$
such that $f_i \in R_0[x_1, \ldots, x_n]$ and such that
$$
S_0 = R_0[x_1, \ldots, x_n]/(f_1, \ldots, f_c)
$$
is a relative global complete intersection.
\end{lemma}

\begin{proof}
Let $R_0 \subset R$ be the $\mathbf{Z}$-algebra of $R$ generated by all the
coefficients of the polynomials $f_1, \ldots, f_c$. Let
$S_0 = R_0[x_1, \ldots, x_n]/(f_1, \ldots, f_c)$.
Clearly, $S = R \otimes_{R_0} S_0$.
Pick a prime $\mathfrak q \subset S$ and denote
$\mathfrak p \subset R$, $\mathfrak q_0 \subset S_0$, and
$\mathfrak p_0 \subset R_0$ the primes it lies over.
Because $\dim (S \otimes_R \kappa(\mathfrak p) ) = n - c$
we also have $\dim (S_0 \otimes_{R_0} \kappa(\mathfrak p_0)) = n - c$,
see Lemma \ref{lemma-dimension-preserved-field-extension}.
By Lemma \ref{lemma-dimension-fibres-bounded-open-upstairs}
there exists a $g \in S_0$, $g \not \in \mathfrak q_0$
such that all nonempty fibres of $R_0 \to (S_0)_g$
have dimension $\leq n - c$. As $\mathfrak q$ was arbitrary and
$\Spec(S)$ quasi-compact, we can find finitely many
$g_1, \ldots, g_m \in S_0$ such that (a) for $j = 1, \ldots, m$
the nonempty fibres of
$R_0 \to (S_0)_{g_j}$ have dimension $\leq n - c$ and (b) the image of
$\Spec(S) \to \Spec(S_0)$ is contained in $D(g_1) \cup \ldots \cup D(g_m)$.
In other words, the images of $g_1, \ldots, g_m$ in $S = R \otimes_{R_0} S_0$
generate the unit ideal. After increasing $R_0$ we may assume
that $g_1, \ldots, g_m$ generate the unit ideal in $S_0$. By (a)
the nonempty fibres of $R_0 \to S_0$ all have dimension $\leq n - c$
and we conclude.
\end{proof}

\begin{lemma}
\label{lemma-relative-global-complete-intersection-conormal}
Let $R$ be a ring. Let $S = R[x_1, \ldots, x_n]/(f_1, \ldots, f_c)$
be a relative global complete intersection (Definition
\ref{definition-relative-global-complete-intersection}). For every prime
$\mathfrak q$ of $S$, let $\mathfrak q'$ denote the corresponding
prime of $R[x_1, \ldots, x_n]$. Then
\begin{enumerate}
\item $f_1, \ldots, f_c$ is a regular sequence in the local ring
$R[x_1, \ldots, x_n]_{\mathfrak q'}$,
\item each of the rings
$R[x_1, \ldots, x_n]_{\mathfrak q'}/(f_1, \ldots, f_i)$ is flat over $R$, and
\item the $S$-module $(f_1, \ldots, f_c)/(f_1, \ldots, f_c)^2$
is free with basis given by the elements $f_i \bmod (f_1, \ldots, f_c)^2$.
\end{enumerate}
\end{lemma}

\begin{proof}
Assume $R$ is Noetherian. Let $\mathfrak p = R \cap \mathfrak q'$.
By Lemma \ref{lemma-lci} for example we see
that $f_1, \ldots, f_c$ form a regular sequence in the local ring
$R[x_1, \ldots, x_n]_{\mathfrak q'} \otimes_R \kappa(\mathfrak p)$.
Moreover, the local ring $R[x_1, \ldots, x_n]_{\mathfrak q'}$
is flat over $R_{\mathfrak p}$. Since $R$, and hence
$R[x_1, \ldots, x_n]_{\mathfrak q'}$ is Noetherian we see from
Lemma \ref{lemma-grothendieck-regular-sequence} that (1) and (2) hold.

\medskip\noindent
Let $R$ be general. Write $R = \colim_{\lambda \in \Lambda} R_\lambda$
as the filtered colimit of finite type $\mathbf{Z}$-subalgebras (compare with
Section \ref{section-colimits-flat}). We may assume that
$f_1, \ldots, f_c \in R_\lambda[x_1, \ldots, x_n]$ for all $\lambda$.
Let $R_0 \subset R$ be as in
Lemma \ref{lemma-relative-global-complete-intersection-Noetherian}.
Then we may assume $R_0 \subset R_\lambda$ for all $\lambda$.
It follows that $S_\lambda = R_\lambda[x_1, \ldots, x_n]/(f_1, \ldots, f_c)$
is a relative global complete intersection (as base change of
$S_0$ via $R_0 \to R_\lambda$, see 
Lemma \ref{lemma-base-change-relative-global-complete-intersection}).
Denote $\mathfrak p_\lambda$, $\mathfrak q_\lambda$, $\mathfrak q'_\lambda$
the prime of $R_\lambda$, $S_\lambda$, $R_\lambda[x_1, \ldots, x_n]$
induced by $\mathfrak p$, $\mathfrak q$, $\mathfrak q'$.
With this notation, we have (1) and (2) for each $\lambda$. Since
$$
R[x_1, \ldots, x_n]_{\mathfrak q'}/(f_1, \ldots, f_i)
=
\colim R_\lambda[x_1, \ldots, x_n]_{\mathfrak q_\lambda'}/(f_1, \ldots, f_i)
$$
we deduce flatness in (2) over $R$ from
Lemma \ref{lemma-colimit-rings-flat}.
Since we have
\begin{align*}
R[x_1, \ldots, x_n]_{\mathfrak q'}/(f_1, \ldots, f_i)
\xrightarrow{f_{i + 1}}
R[x_1, \ldots, x_n]_{\mathfrak q'}/(f_1, \ldots, f_i) \\
=
\colim
\left(
R_\lambda[x_1, \ldots, x_n]_{\mathfrak q_\lambda'}/(f_1, \ldots, f_i)
\xrightarrow{f_{i + 1}}
R_\lambda[x_1, \ldots, x_n]_{\mathfrak q_\lambda'}/(f_1, \ldots, f_i)
\right)
\end{align*}
and since filtered colimits are exact
(Lemma \ref{lemma-directed-colimit-exact})
we conclude that we have (1).

\medskip\noindent
Proof of (3). Denote $N$ the $S$-module
$(f_1, \ldots, f_c)/(f_1, \ldots, f_c)^2$ and $e_i \in N$ the image of $f_i$.
By Lemma \ref{lemma-regular-quasi-regular} and (1) we know that
$e_1, \ldots, e_c$ is a basis of $N_\mathfrak q$ for all
primes $\mathfrak q$ of $S$. By Lemma \ref{lemma-characterize-zero-local}
we conclude that (3) is true.
\end{proof}

\begin{lemma}
\label{lemma-relative-global-complete-intersection}
A relative global complete intersection is syntomic, i.e., flat.
\end{lemma}

\begin{proof}
Let $R \to S$ be a relative global complete intersection.
The fibres are global complete intersections, and
$S$ is of finite presentation over $R$.
Thus the only thing to prove is that $R \to S$ is flat.
This is true by (2) of
Lemma \ref{lemma-relative-global-complete-intersection-conormal}.
\end{proof}

\begin{lemma}
\label{lemma-adjoin-roots}
Suppose that $A$ is a ring, and
$P(x) = x^n + b_1 x^{n-1} + \ldots + b_n \in A[x]$ is
a monic polynomial over $A$. Then there exists a
syntomic, finite free, faithfully flat ring extension
$A \subset A'$ such that $P(x) = \prod_{i = 1, \ldots, n} (x - \beta_i)$
for certain $\beta_i \in A'$.
\end{lemma}

\begin{proof}
Take $A' = A \otimes_R S$, where $R$ and $S$ are as in
Example \ref{example-roots-universal-polynomial},
where $R \to A$ maps $a_i$ to $b_i$, and let
$\beta_i = -1 \otimes \alpha_i$. Observe that $R \to S$ is
faithfully flat and finite free and syntomic by
Lemma \ref{lemma-relative-global-complete-intersection}.
These properties are inherited by the base change
$A \to A'$; some details omitted.
\end{proof}

\begin{lemma}
\label{lemma-syntomic}
Let $R \to S$ be a ring map.
Let $\mathfrak q \subset S$ be a prime lying over
the prime $\mathfrak p$ of $R$.
The following are equivalent:
\begin{enumerate}
\item There exists an element $g \in S$, $g \not \in \mathfrak q$ such that
$R \to S_g$ is syntomic.
\item There exists an element $g \in S$, $g \not \in \mathfrak q$
such that $S_g$ is a relative global complete intersection over $R$.
\item There exists an element $g \in S$, $g \not \in \mathfrak q$,
such that $R \to S_g$ is of finite presentation,
the local ring map $R_{\mathfrak p} \to S_{\mathfrak q}$ is flat, and
the local ring $S_{\mathfrak q}/\mathfrak pS_{\mathfrak q}$ is
a complete intersection ring over $\kappa(\mathfrak p)$ (see
Definition \ref{definition-lci-local-ring}).
\end{enumerate}
\end{lemma}

\begin{proof}
The implication (1) $\Rightarrow$ (3) is Lemma \ref{lemma-lci-at-prime}.
The implication (2) $\Rightarrow$ (1) is
Lemma \ref{lemma-relative-global-complete-intersection}.
It remains to show that (3) implies (2).

\medskip\noindent
Assume (3). After replacing $S$ by $S_g$ for some $g \in S$,
$g\not\in \mathfrak q$ we may assume $S$ is finitely presented over $R$.
Choose a presentation $S = R[x_1, \ldots, x_n]/I$. Let
$\mathfrak q' \subset R[x_1, \ldots, x_n]$ be the prime corresponding
to $\mathfrak q$. Write $\kappa(\mathfrak p) = k$.
Note that $S \otimes_R k = k[x_1, \ldots, x_n]/\overline{I}$ where
$\overline{I} \subset k[x_1, \ldots, x_n]$ is the ideal generated
by the image of $I$. Let $\overline{\mathfrak q}' \subset k[x_1, \ldots, x_n]$
be the prime ideal generated by the image of $\mathfrak q'$.
By Lemma \ref{lemma-lci-at-prime} the equivalent conditions of
Lemma \ref{lemma-lci} hold for $\overline{I}$ and $\overline{\mathfrak q}'$.
Say the dimension of
$\overline{I}_{\overline{\mathfrak q}'}/
\overline{\mathfrak q}'\overline{I}_{\overline{\mathfrak q}'}$
over $\kappa(\overline{\mathfrak q}')$ is $c$.
Pick $f_1, \ldots, f_c \in I$ mapping to a basis of this vector space.
The images $\overline{f}_j \in \overline{I}$ generate
$\overline{I}_{\overline{\mathfrak q}'}$ (by Lemma \ref{lemma-lci}).
Set $S' = R[x_1, \ldots, x_n]/(f_1, \ldots, f_c)$. Let $J$ be the
kernel of the surjection $S' \to S$. Since $S$ is of finite presentation
$J$ is a finitely generated ideal
(Lemma \ref{lemma-compose-finite-type}). Consider the short exact sequence
$$
0 \to J \to S' \to S \to 0
$$
As $S_\mathfrak q$ is flat over $R$ we see that
$J_{\mathfrak q'} \otimes_R k \to S'_{\mathfrak q'} \otimes_R k$
is injective (Lemma \ref{lemma-flat-tor-zero}).
However, by construction $S'_{\mathfrak q'} \otimes_R k$
maps isomorphically to $S_\mathfrak q \otimes_R k$. Hence we
conclude that $J_{\mathfrak q'} \otimes_R k =
J_{\mathfrak q'}/\mathfrak pJ_{\mathfrak q'} = 0$. By Nakayama's
lemma (Lemma \ref{lemma-NAK}) we conclude that there exists a
$g \in R[x_1, \ldots, x_n]$, $g \not \in \mathfrak q'$ such that
$J_g = 0$. In other words $S'_g \cong S_g$. After further localizing
we see that $S'$ (and hence $S$) becomes a relative global complete
intersection by
Lemma \ref{lemma-localize-relative-complete-intersection}
as desired.
\end{proof}

\begin{lemma}
\label{lemma-syntomic-presentation-ideal-mod-squares}
Let $R$ be a ring. Let $S = R[x_1, \ldots, x_n]/I$ for some
finitely generated ideal $I$. If $g \in S$ is such that
$S_g$ is syntomic over $R$, then $(I/I^2)_g$ is a finite projective
$S_g$-module.
\end{lemma}

\begin{proof}
By Lemma \ref{lemma-syntomic} there exist finitely many elements
$g_1, \ldots, g_m \in S$ which generate the unit ideal in $S_g$
such that each $S_{gg_j}$ is a relative global complete intersection
over $R$. Since it suffices to prove that $(I/I^2)_{gg_j}$ is
finite projective, see
Lemma \ref{lemma-finite-projective},
we may assume that $S_g$ is a relative global complete intersection.
In this case the result follows from
Lemmas \ref{lemma-conormal-module-localize} and
\ref{lemma-relative-global-complete-intersection-conormal}.
\end{proof}

\begin{lemma}
\label{lemma-composition-syntomic}
Let $R \to S$, $S \to S'$ be ring maps.
\begin{enumerate}
\item If $R \to S$ and $S \to S'$ are syntomic, then $R \to S'$
is syntomic.
\item If $R \to S$ and $S \to S'$ are relative global complete intersections,
then $R \to S'$ is a relative global complete intersection.
\end{enumerate}
\end{lemma}

\begin{proof}
Proof of (2). Say $R \to S$ and $S \to S'$ are relative global complete
intersections and we have presentations
$S =  R[x_1, \ldots, x_n]/(f_1, \ldots, f_c)$ and
$S' = S[y_1, \ldots, y_m]/(h_1, \ldots, h_d)$ as in
Definition \ref{definition-relative-global-complete-intersection}.
Then
$$
S' \cong
R[x_1, \ldots, x_n, y_1, \ldots, y_m]/(f_1, \ldots, f_c, h'_1, \ldots, h'_d)
$$
for some lifts $h_j' \in R[x_1, \ldots, x_n, y_1, \ldots, y_m]$ of the $h_j$.
Hence it suffices to bound the dimensions of the fibre rings.
Thus we may assume $R = k$ is a field.
In this case we see that we have a ring, namely $S$, which is of finite
type over $k$ and equidimensional of dimension $n - c$, and a
finite type ring map $S \to S'$ all of whose nonempty fibre
rings are equidimensional of dimension $m - d$. Then, by
Lemma \ref{lemma-dimension-base-fibre-total} for example applied
to localizations at maximal ideals of $S'$, we see that
$\dim(S') \leq n - c + m - d$ as desired.

\medskip\noindent
We will reduce part (1) to part (2). Assume $R \to S$ and $S \to S'$
are syntomic. Let $\mathfrak q' \subset S$ be a prime ideal lying
over $\mathfrak q \subset S$. By Lemma \ref{lemma-syntomic}
there exists a $g' \in S'$, $g' \not \in \mathfrak q'$ such that
$S \to S'_{g'}$ is a relative global complete intersection.
Similarly, we find $g \in S$, $g \not \in \mathfrak q$ such that
$R \to S_g$ is a relative global complete intersection.
By Lemma \ref{lemma-base-change-relative-global-complete-intersection}
the ring map $S_g \to S_{gg'}$ is a relative global complete intersection.
By part (2) we see that $R \to S_{gg'}$ is a relative global
complete intersection and $gg' \not \in \mathfrak q'$.
Since $\mathfrak q'$ was arbitrary
combining Lemmas \ref{lemma-syntomic} and \ref{lemma-local-syntomic}
we see that $R \to S'$ is syntomic (this also uses that the spectrum
of $S'$ is quasi-compact, see Lemma \ref{lemma-quasi-compact}).
\end{proof}

\noindent
The following lemma will be improved later, see
Smoothing Ring Maps, Proposition \ref{smoothing-proposition-lift-smooth}.

\begin{lemma}
\label{lemma-lift-syntomic}
Let $R$ be a ring and let $I \subset R$ be an ideal.
Let $R/I \to \overline{S}$ be a syntomic map.
Then there exists elements $\overline{g}_i \in \overline{S}$
which generate the unit ideal of $\overline{S}$
such that each $\overline{S}_{g_i} \cong S_i/IS_i$
for some relative global complete intersection $S_i$
over $R$.
\end{lemma}

\begin{proof}
By Lemma \ref{lemma-syntomic} we find a collection of elements
$\overline{g}_i \in \overline{S}$
which generate the unit ideal of $\overline{S}$
such that each $\overline{S}_{g_i}$ is a relative
global complete intersection over $R/I$.
Hence we may assume that $\overline{S}$ is a
relative global complete intersection.
Write
$\overline{S} =
(R/I)[x_1, \ldots, x_n]/(\overline{f}_1, \ldots, \overline{f}_c)$
as in Definition \ref{definition-relative-global-complete-intersection}.
Choose $f_1, \ldots, f_c \in R[x_1, \ldots, x_n]$
lifting $\overline{f}_1, \ldots, \overline{f}_c$.
Set $S = R[x_1, \ldots, x_n]/(f_1, \ldots, f_c)$.
Note that $S/IS \cong \overline{S}$.
By Lemma \ref{lemma-localize-relative-complete-intersection}
we can find $g \in S$ mapping to $1$ in $\overline{S}$ such
that $S_g$ is a relative global complete intersection over $R$.
Since $\overline{S} \cong S_g/IS_g$ this finishes the proof.
\end{proof}


